\chapterrule
\section{从三维到真实模型：规模、配置与分词}
\label{intro:23a-from-toy-to-gpt}
\chaptercompass{三维模型与 GPT-2 的差别是什么？}{训练任务相似，对照真实模型的规模、Tokenizer 与参数细节；数据工程见训练系统章。}{$\mathrm{IDs}:1\times T\to Z:T\times|\mathcal V|$}{GPT-2 的 Embedding、投影、Norm 等参数；本书三维主例只有 99 个训练元素}
\modelpipelinefocus{7}{真实模型续写推理}


上一章会用训练好的模型一步步生成 Token。
但我们手算的网络只有一层、一个头、三个隐藏特征和五项词表。
它怎样展示 GPT 的训练原理，真实模型还要扩大哪些部分？

\subsection{回看先前的 GPT-2 small 尺度示例}

前面的\hyperref[intro:20b-task-contract]{语言预测任务}已给出尺寸对照。
按 \href{https://huggingface.co/openai-community/gpt2/blob/607a30d783dfa663caf39e06633721c8d4cfcd7e/config.json}{GPT-2 small 的公开配置}：
词表 $50{,}257$，隐藏宽度 $768$，
$12$ 层 Transformer Block、$12$ 个注意力头，最多处理 $1{,}024$ 个位置；
每层 FFN 的中间宽度是隐藏宽度的 $4$ 倍，即 $4\times768=3{,}072$。
后来的模型普遍更宽、更深，并会替换其中的部件（见第~\ref*{intro:22d-modern-transformers}~章），
但都可以按这几项来读配置。

\subsection{宽度、层数、头数分别在改变什么？}

\textbf{宽度}增加每个 Token 的隐藏特征维度。
本书主例 $d_{\rm model}=3$，所有隐藏行都是 $1\times3$；
GPT-2 small 为 $1\times768$。

\textbf{头数}让不同注意力头在不同投影子空间中读取内容。
若将 768 个隐藏维度等分给 12 个注意力头，
每头 Query/Key 的宽度为 $768/12=64$。
多头要求 $d_{\rm model}$ 能被头数整除；为让每头至少有 $2$ 维，
双头例另用 $d_{\rm model}=4$，三维主例保留单头。

\textbf{层数}允许把一个 Block 的输出作为下一个 Block 的输入，
让上下文特征反复组合。它不是简单把每层的预测概率相加。

\textbf{FFN 宽度}决定每个 Token 经过非线性中间层时可用的特征数量，
不是 Token 的个数；它与词表输出头也不是同一个矩阵。
GPT-2 的 FFN 激活函数是 GELU（配置中的 \texttt{gelu\_new}），
本书主例用 ReLU；两者都是逐元素非线性，GELU 在零附近平滑过渡，负输入给出接近零的小负值。

\subsection{参数怎么数：99 个教学参数与 GPT-2 量级}

本书三维贯通模型\textbf{按实际可训练元素个数}计数。
固定 LayerNorm 仿射系数、不使用偏置，输出头与输入 Embedding 独立，
也不把计算图中的中间激活当作参数：
\[
\begin{aligned}
E &: 5\times3=15, & P_{\rm pos} &:3\times3=9,\\
W_Q,W_K,W_V,W_O &:4\times(3\times3)=36,\\
W_{\rm up},W_{\rm down} &:3\times4+4\times3=24,\\
W_{\rm out} &:3\times5=15.
\end{aligned}
\]
因此总数为
\begin{equation}
\boxed{15+9+36+24+15=99\text{ 个可训练参数}.}
\label{eq:toy-99-parameters}
\end{equation}
\textbf{运行时显存还包含额外状态}：
推理还会保存激活、缓存和程序开销，训练还需梯度及优化器状态。

再用经典 GPT-2 small 的已公开结构做量级估算：
Token Embedding 约有
$50{,}257\times768=38{,}597{,}376$ 个参数，约 $38.6$ M；
每个 Block 的四组 Attention 投影与两个约 $4d$ 的 FFN 投影，
主导项约为 $12d^2$，
即 $12\times768^2\approx7.08$ M/层。
十二层约 $84.9$ M；再加可学习位置 Embedding
$1{,}024\times768\approx0.79$ M、
归一化参数与偏置，得到约 $124$ M 的经典量级。
这是\textbf{结构上的近似预算}；逐项计入偏置与归一化参数，精确值为 $124{,}439{,}808$。
GPT-2 的输出头复用 Token Embedding，不再另加约 $38.6$ M；
若不共享，总数约为 $163$ M（精确为 $163{,}037{,}184$）。

\subsection{［复算］标准 GPT Block 的 FFN 为什么比 Attention 参数多？}

在\textbf{经典 Dense Transformer} 的简化预算中，
注意力有 $W_Q,W_K,W_V,W_O$ 四个 $d\times d$ 投影；
两层普通 FFN 则先从 $d$ 扩到 $4d$，再降回 $d$：
\[
\underbrace{N_{\rm att}}_{4d^2}=4d^2,\qquad
\underbrace{N_{\rm FFN}}_{8d^2}=d(4d)+(4d)d=8d^2.
\]
因此在\textbf{不计偏置与归一化参数}的前提下，
一层总主导项为 $12d^2$，其中 FFN 约占三分之二。

\begin{center}\small
\tagpdfsetup{table/header-rows={1}}
\begin{tabularx}{\linewidth}{@{}p{24mm}p{28mm}X@{}}
\toprule
部件 & 主导公式 & $d=768$ 时参数元素数\\\midrule
Q/K/V/O & $4d^2$ & $2{,}359{,}296$\\
Dense FFN & $8d^2$ & $4{,}718{,}592$\\
两者合计 & $12d^2$ & $7{,}077{,}888$\\
\bottomrule
\end{tabularx}
\end{center}
这些是\textbf{参数数量}，不是 FLOPs、KV Cache 或实际速度。
SwiGLU、MoE、共享投影和不同 FFN 扩展比例都会改变这个对照；
其他 FFN 宽度和门控结构要重新按矩阵尺寸计数。


\subsection{［复算］训练前损失应该与什么比较？}

如果词表有 $|\mathcal V|$ 个候选且每个 Token 的概率\textbf{恰好均匀}，
那么任一正确类别的交叉熵为
\[
-\ln\frac1{|\mathcal V|}=\ln|\mathcal V|.
\]
本书五项词表的基准是
$\ln5\approx\boxed{1.609}$；
GPT-2 small 的词表规模为 $50{,}257$，
对应均匀基准 $\ln 50{,}257\approx\boxed{10.825}$。

\textbf{均匀概率给出对照基准；实际初始损失要按模型输出计算}：
实际 logits 的分散程度、初始化与数据分布都会影响初始交叉熵。
本书三维例子的初始损失约为 $1.798430$，
高于均匀基准并不矛盾，因为我们手工设定了非均匀的初始参数和输出概率。
先确认基线量级、再尝试\textbf{过拟合同一小批数据}，
是检查训练链路是否正常的常见调试方法。

\subsection{训练目标相同，为什么模型能力仍有差距？}

无论玩具模型还是 GPT-2 类 Decoder-only 模型，
在自回归预训练时都可以从当前及先前 Token 预测下一 Token。
不同点在于训练文本规模、Tokenization、
模型宽度和深度、参数量、优化预算以及数据分布。
本书的单条训练序列从 $L=1.7984$ 降到 $1.6114$，
验证了\textbf{这组参数和数据上的一次优化}。
真实语言理解、事实核验及指令遵循要靠独立任务评测，
而不是观察单次训练损失。

原始 GPT-2 的输出头与输入 Token Embedding 权重\textbf{绑定}，
采用词嵌入矩阵的转置来得到词表分数。
本书为了让每步梯度和形状易于检查，采用\textbf{独立的 $3\times5$ 输出头}，
这两种输出头的参数计数口径不同；
真实模型还要计入位置 Embedding、归一化和偏置等部分。
其他模型是否绑定权重，应检查其配置与输出头实现。

\subsection{\texorpdfstring{{\symboliconfont ★}}{★}选读复算：真实 GPT-2 的分词、概率、温度与注意力}

\localfixture{公开 GPT-2 small 预训练权重；非本书三维参数}
这一节提供一个\textbf{独立运行的真实模型实验}：
\href{https://huggingface.co/openai-community/gpt2}{openai-community/gpt2}，
固定修订（短 SHA）
\texttt{607a30d7}；完整 40 位 revision
由 \texttt{scripts/gpt2\_demo.py} 固定并在运行时输出。
GPT-2 small 主要在英文语料上预训练，
适合观察真实 Token、词表分数及 Attention，
本节用它观察英文续写的计算过程，下一小节再对照中文分词。

先在仓库根目录运行无需联网或额外依赖的数学自检：
\begin{verbatim}
python3 scripts/gpt2_demo.py
\end{verbatim}
真实推理需要先自行安装兼容的 \texttt{torch}、
\texttt{transformers}，绘制热图还需 \texttt{matplotlib}；
模型文件首次运行时需下载。显式启用：
\begin{verbatim}
python3 scripts/gpt2_demo.py --real
\end{verbatim}
已有缓存可加 \texttt{--local-files-only}；\texttt{--heatmap gpt2-attn.png} 导出首层首头图。
脚本核对四项：真实 Token 数、$1\times T\times50{,}257$ logits、各温度概率和、未来位置权重为零。
完整环境、后端与运行记录集中见附录 G。

\begin{example}
\noindent\emph{一份真实 GPT-2 small 的复算记录}\par
在公开权重修订 \texttt{607a30d7}、CPU PyTorch 2.5.1、
Transformers 4.46.3 下，
英文 Prompt 为 \texttt{The capital of France is}。
这次真实推理得到：

\begin{center}\small
\tagpdfsetup{table/header-rows={1}}
\begin{tabularx}{\linewidth}{@{}p{41mm}X@{}}
\toprule
实测项目 & 输出或检查结果\\\midrule
真实 Token 数与 ID & $5$；$(464,3139,286,4881,318)$\\
全位置 logits & $1\times5\times50{,}257$\\
温度 $0.7,1.0,1.5$ & 各温度的概率总和均为 $1.00000000$\\
首层首头未来位置 & Attention 上三角最大绝对值为 $0$\\
\bottomrule
\end{tabularx}
\end{center}

温度 $1.0$ 时，下一 Token 的最高概率候选是
\texttt{the}（含前导空格，ID 262，约 $0.08459$）；
\texttt{Paris}（含前导空格，ID 6342）约为 $0.03224$。
下一 Token 分布回答的是怎样续写；事实回答需要外部出处支撑。


本次实验的固定依赖、Top-5 与注意力热图复现入口统一放在
\hyperref[app:revision-notes]{附录 G 的真实 GPT-2 运行记录}。
上述实测值仅对应所述输入、权重及运行环境。
\end{example}

GPT-2 属于 Decoder-only 家族；与原论文翻译架构的对照见
\sectionentry{elective:original-paper}。

\subsection{\texorpdfstring{{\symboliconfont ★}}{★}选读复算：同一句中文，真实分词器怎样编号？}
\phantomsection\label{experiment:real-chinese-tokenizer}

回到“春天花开了。”。沿用第~\ref*{intro:20a-tokenizer}~章锁定的
Qwen2.5-0.5B 分词器，\texttt{add\_special\_tokens=False} 时实测为：
\begin{center}
\tagpdfsetup{table/header-rows={1}}
\begin{tabular}{lrr}
\toprule
解码片段 & 真实 Token ID & 教学词表 ID\\\midrule
春天 & 105303 & 2\\
花 & 99232 & 4\\
开了 & 101467 & 5\\
。 & 1773 & 3\\
\bottomrule
\end{tabular}
\end{center}
带句号共四个 Token，去掉句号的“春天花开了”共三个。
\textbf{这句话的切分恰好与教学切分一致，编号却来自不同词表}。
真实 ID 直接用作分词器所属模型的数组索引；本书教学 ID 从 1 起，写 Python 时才减 1。
输入别的句子应重新运行分词器，不能把这四个片段推成通用中文分词规则。

% Keep this closing block with the self-check so the part-end page is not a sparse tail.
\Needspace{12\baselineskip}
运行：
\begin{verbatim}
python3 scripts/chinese_tokenizer_demo.py --real
\end{verbatim}
完整权重修订\par\nobreak
\noindent\texttt{060db6499f32faf8b98477b0a26969ef7d8b9987}、
分词器文件哈希和实测记录见仓库 \path{scripts/fixtures/qwen25_tokenizer_trace.json}。
训练资料的来源、去重与配比则转到第~\ref*{intro:23b-deployment-and-ecosystem}~章的
\sectionentry{sec:data-engineering}，分别检查模型与数据两条链。

\chapterchecktasks{23a-from-toy-to-gpt}{正文基础三题}{
\item \textbf{形状：}GPT-2 small 的隐藏宽度为 $768$，采用 $12$ 个等宽头。每个头宽是多少？这是否直接决定模型拥有多少层？\chapterinlinehint{用隐藏宽度除以头数，再单独检查层数配置。}
\item \textbf{估算：}忽略偏置，令每层 Attention 投影参数约 $4d^2$、FFN 参数约 $8d^2$。代入 $d=768$，哪个组件占更多层内参数？\chapterinlinehint{先比较两个参数公式的系数。}
\item \textbf{找错：}同学在三维玩具任务上降低了损失，便声称它已具备大型模型的能力。还缺哪两类独立证据？\chapterinlinehint{分别列出未见任务评测与事实核验所需证据。}
}
\chapterrecallanswer{玩具与 GPT 可使用相同计算框架，但宽度、层数、数据、参数与外部评测共同决定模型能力。}





\begin{conclusion}
从三维模型到真实预训练，计算链相通；规模、数据质量与独立评测共同决定效果。
下一部分把模型接到证据、工具与授权边界；图像输入作为扩展路径，放在\hyperref[part:vision]{视觉部分}。
\end{conclusion}
